\chapter{Stokes' Theorem}
\label{ch:vii11-stokes-theorem}

Stokes' theorem is the global statement that exterior differentiation and oriented boundary are dual operations:
\[
\boxed{\int_M d\omega=\int_{\partial M}\omega.}
\]
It contains the fundamental theorem of calculus, Green's theorem, the classical Stokes theorem in $\mathbb R^3$, and the divergence theorem as coordinate expressions of one geometric identity.

\section{Manifolds with boundary}
\begin{definition}[Smooth manifold with boundary]\label{def:vii11-boundary-manifold}
An $n$-manifold with boundary is locally modeled on the half-space
\[
\mathbb H^n=\{x\in\mathbb R^n:x^n\ge0\}.
\]
Points mapping to $x^n=0$ form the boundary $\partial M$; the remaining points form the interior $M^\circ$.
\end{definition}

\begin{theorem}[Boundary is a manifold]\label{thm:vii11-boundary-manifold}
If $M$ is a smooth $n$-manifold with boundary, then $\partial M$ is a smooth $(n-1)$-manifold without boundary.
\end{theorem}

\begin{definition}[Outward-normal-first orientation]\label{def:vii11-boundary-orientation}
Let $M$ be oriented. A basis $(v_1,\ldots,v_{n-1})$ of $T_p\partial M$ is positive when $(\nu,v_1,\ldots,v_{n-1})$ is a positive basis of $T_pM$ for any outward-pointing transverse vector $\nu$.
\end{definition}

\begin{example}[Interval boundary signs]\label{ex:vii11-interval}
For the standard orientation of $[a,b]$, the induced orientation on its $0$-dimensional boundary is
\[
\partial[a,b]=\{b\}-\{a\}.
\]
\end{example}

\section{Stokes' theorem}
\begin{theorem}[Stokes]\label{thm:vii11-stokes}
Let $M$ be an oriented smooth $n$-manifold with boundary and let $\omega\in\Omega_c^{n-1}(M)$. Then
\[
\int_M d\omega=\int_{\partial M}\iota^*\omega,
\]
where $\iota:\partial M\hookrightarrow M$ is inclusion and the boundary carries the outward-normal-first orientation.
\end{theorem}

\begin{remark}[Boundaryless case]\label{rem:vii11-boundaryless}
If $\partial M=\varnothing$, then every compactly supported exact top form has integral zero:
\[
\int_Md\omega=0.
\]
\end{remark}

\section{Why the theorem is true}
\begin{proposition}[Local half-space calculation]\label{prop:vii11-local-halfspace}
For a compactly supported $(n-1)$-form on $\mathbb H^n$, Stokes reduces to the one-variable fundamental theorem of calculus in the normal coordinate.
\end{proposition}
\begin{proof}
Write $\omega=\sum_i(-1)^{i-1}f_i\,dx^1\wedge\cdots\widehat{dx^i}\cdots\wedge dx^n$. Then
\[
d\omega=\left(\sum_i\partial_i f_i\right)dx^1\wedge\cdots\wedge dx^n.
\]
All tangential derivative integrals vanish by compact support. The $x^n$ derivative contributes the boundary value at $x^n=0$ with the sign fixed by the outward orientation.
\end{proof}

\begin{proof}[Proof of Stokes' theorem]
Choose a partition of unity subordinate to interior charts and boundary half-space charts. Apply the Euclidean compact-support calculation to each localized form. Interior pieces have no boundary contribution; boundary pieces contribute exactly the pullback to $\partial M$. Summing gives the theorem.
\end{proof}

\section{The fundamental theorem of calculus}
\begin{corollary}[Fundamental theorem]\label{cor:vii11-ftc}
For a smooth function $f$ on $[a,b]$,
\[
\int_{[a,b]}df=f(b)-f(a).
\]
\end{corollary}

\section{Green's theorem}
\begin{corollary}[Green]\label{cor:vii11-green}
For a positively oriented planar region $D$ with smooth boundary and smooth functions $P,Q$,
\[
\oint_{\partial D}P\,dx+Q\,dy
=\iint_D\left(\frac{\partial Q}{\partial x}-\frac{\partial P}{\partial y}\right)dx\,dy.
\]
\end{corollary}
\begin{proof}
Apply Stokes to $\omega=P\,dx+Q\,dy$, for which
\[
d\omega=(Q_x-P_y)\,dx\wedge dy.
\]
\end{proof}

\section{Classical Stokes in three dimensions}
\begin{corollary}[Curl theorem]\label{cor:vii11-curl}
Let $S\subset\mathbb R^3$ be an oriented smooth surface with boundary and $F=(P,Q,R)$. For
\[
\omega=P\,dx+Q\,dy+R\,dz,
\]
Stokes' theorem becomes
\[
\int_{\partial S}F\cdot d\mathbf r
=\int_S(\nabla\times F)\cdot n\,dS.
\]
\end{corollary}

\section{Divergence theorem}
\begin{corollary}[Divergence theorem]\label{cor:vii11-divergence}
For a compact domain $M\subset\mathbb R^n$ with smooth boundary and a vector field $X=\sum_iX^i\partial_i$,
\[
\int_M\operatorname{div}X\,dV=\int_{\partial M}X\cdot\nu\,dS.
\]
\end{corollary}
\begin{proof}
Let $dV=dx^1\wedge\cdots\wedge dx^n$ and set $\omega=\iota_XdV$. Then
\[
d\omega=(\operatorname{div}X)dV,
\]
and the boundary pullback is the flux form.
\end{proof}

\section{Consequences for exact and closed forms}
\begin{definition}[Closed and exact]\label{def:vii11-closed-exact}
A differential form $\alpha$ is closed if $d\alpha=0$ and exact if $\alpha=d\beta$ for some form $\beta$. Since $d^2=0$, every exact form is closed.
\end{definition}

\begin{corollary}[Exact forms on cycles]\label{cor:vii11-exact-cycle}
If $C$ is a compact oriented $k$-manifold without boundary and $\alpha=d\beta$ is an exact $k$-form defined near $C$, then
\[
\int_C\alpha=0.
\]
\end{corollary}

\begin{corollary}[Boundary of a boundary]\label{cor:vii11-boundary-boundary}
For an oriented manifold with corners treated facewise, codimension-two boundary contributions cancel in pairs, reflecting the algebraic identity $\partial^2=0$ and the differential identity $d^2=0$.
\end{corollary}

\section{Boundary orientation checks}
\begin{example}[Unit disk]\label{ex:vii11-disk}
The standard orientation $dx\wedge dy$ on the unit disk induces the counterclockwise orientation on $S^1$.
\end{example}

\begin{example}[Upper half-plane]\label{ex:vii11-halfplane}
For $\mathbb H^2=\{(x,y):y\ge0\}$ with orientation $(\partial_x,\partial_y)$, the outward vector along $y=0$ is $-\partial_y$. Therefore the positive boundary tangent is $\partial_x$.
\end{example}


\section{Exercises}

\begin{exercise}\label{exr:vii11-01}
Identify the boundary and interior of $[0,1]$ as a $1$-manifold with boundary.
\end{exercise}
\begin{hint}
Use the half-line model.
\end{hint}
\begin{solution}
The boundary is $\{0,1\}$ and the interior is $(0,1)$.
\end{solution}

\begin{exercise}\label{exr:vii11-02}
Show that $\partial M$ is locally modeled on $\mathbb R^{n-1}$.
\end{exercise}
\begin{hint}
Restrict a boundary chart to $x^n=0$.
\end{hint}
\begin{solution}
A half-space chart restricts to a homeomorphism from the boundary slice onto an open subset of $\mathbb R^{n-1}$, and transition maps restrict smoothly.
\end{solution}

\begin{exercise}\label{exr:vii11-03}
For the standard orientation of $[a,b]$, compute the induced signs of the endpoints.
\end{exercise}
\begin{hint}
Use outward-normal-first.
\end{hint}
\begin{solution}
At $b$ the outward vector agrees with the positive tangent, giving sign $+1$; at $a$ it is negative, giving sign $-1$.
\end{solution}

\begin{exercise}\label{exr:vii11-04}
Show that the standard orientation on the unit disk induces counterclockwise orientation on its boundary.
\end{exercise}
\begin{hint}
Test at $(1,0)$.
\end{hint}
\begin{solution}
At $(1,0)$ the outward normal is $\partial_x$; the boundary vector $\partial_y$ makes $(\partial_x,\partial_y)$ positive, so the induced direction is counterclockwise.
\end{solution}

\begin{exercise}\label{exr:vii11-05}
State Stokes' theorem for a compactly supported $(n-1)$-form.
\end{exercise}
\begin{hint}
Include the boundary pullback.
\end{hint}
\begin{solution}
$\int_Md\omega=\int_{\partial M}\iota^*\omega$ with the outward-normal-first boundary orientation.
\end{solution}

\begin{exercise}\label{exr:vii11-06}
Derive the fundamental theorem of calculus from Stokes.
\end{exercise}
\begin{hint}
Take a $0$-form $f$ on an interval.
\end{hint}
\begin{solution}
Stokes gives $\int_{[a,b]}df=\int_{\{b\}-\{a\}}f=f(b)-f(a)$.
\end{solution}

\begin{exercise}\label{exr:vii11-07}
Compute $d(P\,dx+Q\,dy)$.
\end{exercise}
\begin{hint}
Use $d(f\alpha)=df\wedge\alpha$.
\end{hint}
\begin{solution}
$d(P\,dx+Q\,dy)=(Q_x-P_y)\,dx\wedge dy$.
\end{solution}

\begin{exercise}\label{exr:vii11-08}
Recover Green's theorem from Stokes.
\end{exercise}
\begin{hint}
Use the preceding one-form.
\end{hint}
\begin{solution}
Substitute $d\omega=(Q_x-P_y)dx\wedge dy$ into Stokes on the oriented region.
\end{solution}

\begin{exercise}\label{exr:vii11-09}
Show that $d^2f=0$ for a smooth function.
\end{exercise}
\begin{hint}
Use equality of mixed partials.
\end{hint}
\begin{solution}
In coordinates the coefficient of $dx^i\wedge dx^j$ is $\partial_i\partial_jf-\partial_j\partial_if=0$.
\end{solution}

\begin{exercise}\label{exr:vii11-10}
Why is every exact form closed?
\end{exercise}
\begin{hint}
Apply $d$ twice.
\end{hint}
\begin{solution}
If $\alpha=d\beta$, then $d\alpha=d^2\beta=0$.
\end{solution}

\begin{exercise}\label{exr:vii11-11}
Show that an exact top form integrates to zero on a compact oriented manifold without boundary.
\end{exercise}
\begin{hint}
Apply Stokes.
\end{hint}
\begin{solution}
If $\alpha=d\beta$, then $\int_M\alpha=\int_{\partial M}\beta=0$.
\end{solution}

\begin{exercise}\label{exr:vii11-12}
Compute $\oint_{S^1}(-y\,dx+x\,dy)$ by Stokes.
\end{exercise}
\begin{hint}
Differentiate the one-form over the disk.
\end{hint}
\begin{solution}
Its differential is $2\,dx\wedge dy$, so the integral is $2$ times the unit-disk area, namely $2\pi$.
\end{solution}

\begin{exercise}\label{exr:vii11-13}
For $F=(P,Q,R)$, relate $d(Pdx+Qdy+Rdz)$ to curl.
\end{exercise}
\begin{hint}
Expand the exterior derivative.
\end{hint}
\begin{solution}
The coefficients of $dy\wedge dz$, $dz\wedge dx$, $dx\wedge dy$ are precisely the components of $\nabla\times F$.
\end{solution}

\begin{exercise}\label{exr:vii11-14}
Use Stokes to show the circulation of a gradient around a closed curve is zero.
\end{exercise}
\begin{hint}
The corresponding one-form is exact.
\end{hint}
\begin{solution}
If $F=\nabla f$, then $F\cdot d\mathbf r=df$, and the integral over a closed one-manifold is zero.
\end{solution}

\begin{exercise}\label{exr:vii11-15}
For $X=\sum_iX^i\partial_i$, compute $d(\iota_XdV)$ in $\mathbb R^n$.
\end{exercise}
\begin{hint}
Differentiate each coefficient.
\end{hint}
\begin{solution}
All cross terms vanish because of repeated differentials, leaving $(\sum_i\partial_iX^i)dV=(\operatorname{div}X)dV$.
\end{solution}

\begin{exercise}\label{exr:vii11-16}
Recover the divergence theorem from Stokes.
\end{exercise}
\begin{hint}
Use $\omega=\iota_XdV$.
\end{hint}
\begin{solution}
Stokes gives $\int_M(\operatorname{div}X)dV=\int_{\partial M}\iota^*(\iota_XdV)$, and the boundary form is the outward flux $X\cdot\nu\,dS$.
\end{solution}

\begin{exercise}\label{exr:vii11-17}
Show that reversing the orientation of $M$ reverses both sides of Stokes.
\end{exercise}
\begin{hint}
Track induced boundary orientation too.
\end{hint}
\begin{solution}
The interior integral changes sign, and the boundary orientation induced by outward-normal-first also reverses, so the boundary integral changes sign as well.
\end{solution}

\begin{exercise}\label{exr:vii11-18}
What happens to the induced boundary orientation if the outward vector is replaced by another outward transverse vector?
\end{exercise}
\begin{hint}
Connect the vectors through the outward half-space.
\end{hint}
\begin{solution}
Any two outward transverse vectors differ by a positive normal component plus tangential terms; the change-of-basis determinant is positive, so the induced orientation is unchanged.
\end{solution}

\begin{exercise}\label{exr:vii11-19}
Use Stokes to show there is no $1$-form $\beta$ on $S^2$ with $d\beta$ equal to the standard area form.
\end{exercise}
\begin{hint}
Integrate over $S^2$.
\end{hint}
\begin{solution}
If the area form were exact, its integral would be zero by Stokes, but its integral is $4\pi$.
\end{solution}

\begin{exercise}\label{exr:vii11-20}
For the annulus with standard orientation, determine the boundary directions.
\end{exercise}
\begin{hint}
Outward normals point opposite ways on the two circles.
\end{hint}
\begin{solution}
The outer circle is counterclockwise; the inner circle is clockwise.
\end{solution}

\begin{exercise}\label{exr:vii11-21}
Why do interior partition-of-unity pieces contribute no boundary term in the proof of Stokes?
\end{exercise}
\begin{hint}
Their support avoids the boundary.
\end{hint}
\begin{solution}
Such a localized form is compactly supported in an open Euclidean chart, where integration of a total derivative vanishes by the ordinary fundamental theorem and compact support.
\end{solution}

\begin{exercise}\label{exr:vii11-22}
Explain why compact support is needed in the noncompact version of Stokes.
\end{exercise}
\begin{hint}
Prevent terms at infinity.
\end{hint}
\begin{solution}
Compact support ensures all coordinate integrations by parts have no contribution at infinity and makes the partition sum finite near the support.
\end{solution}

\begin{exercise}\label{exr:vii11-23}
Show that $\int_{\partial M}\iota^*d\eta=0$ when $\partial M$ is compact and has no boundary.
\end{exercise}
\begin{hint}
Apply Stokes on $\partial M$.
\end{hint}
\begin{solution}
$\int_{\partial M}d(\iota^*\eta)=\int_{\partial(\partial M)}\iota^*\eta=0$.
\end{solution}

\begin{exercise}\label{exr:vii11-24}
Explain geometrically why $d^2=0$ mirrors $\partial^2=0$.
\end{exercise}
\begin{hint}
Pair Stokes twice.
\end{hint}
\begin{solution}
For suitable $\eta$, $\int_Md^2\eta=\int_{\partial M}d\eta=\int_{\partial^2M}\eta$. The algebraic zero on the left corresponds to cancellation of codimension-two boundary faces on the right.
\end{solution}


\section{Solved problem dossiers}

\begin{problem}[Legacy Stokes problem: interval]\label{prob:vii11-01}
Derive the fundamental theorem of calculus from the oriented-boundary convention.
\end{problem}
\begin{solution}
The standard orientation on $[a,b]$ gives $\partial[a,b]=\{b\}-\{a\}$. Applying Stokes to the $0$-form $f$ gives $\int_a^bf'(t)dt=f(b)-f(a)$.
\end{solution}

\begin{problem}[Legacy Stokes problem: planar Green theorem]\label{prob:vii11-02}
Derive Green's theorem from the exterior derivative of $Pdx+Qdy$.
\end{problem}
\begin{solution}
Compute $d(Pdx+Qdy)=(Q_x-P_y)dx\wedge dy$, then apply Stokes with the counterclockwise induced boundary orientation.
\end{solution}

\begin{problem}[Legacy Stokes problem: boundary orientation]\label{prob:vii11-03}
Prove the unit disk induces counterclockwise orientation on the unit circle.
\end{problem}
\begin{solution}
At each point choose the outward radial vector first. The tangent obtained by rotating it counterclockwise by $\pi/2$ completes it to the standard positive basis of the plane.
\end{solution}

\begin{problem}[Local half-space proof]\label{prob:vii11-04}
Prove Stokes for a compactly supported $(n-1)$-form on $\mathbb H^n$.
\end{problem}
\begin{solution}
Expand into coordinate omissions. Tangential derivatives integrate to zero by compact support. The normal derivative integrates by the one-dimensional fundamental theorem to the boundary value, with the outward-normal-first sign.
\end{solution}

\begin{problem}[Partition-of-unity proof]\label{prob:vii11-05}
Globalize the half-space calculation to an oriented manifold with boundary.
\end{problem}
\begin{solution}
Choose a subordinate partition of unity. Apply the interior Euclidean and boundary half-space calculations to each localized form. Sum; the boundary restrictions of the partition still sum to one.
\end{solution}

\begin{problem}[Classical curl theorem]\label{prob:vii11-06}
Translate Stokes for a one-form in $\mathbb R^3$ into the usual line-integral/curl formula.
\end{problem}
\begin{solution}
Identify $Pdx+Qdy+Rdz$ with the vector field $F$. Its exterior derivative corresponds to $\nabla\times F$, while pullback to the boundary curve gives $F\cdot d\mathbf r$.
\end{solution}

\begin{problem}[Divergence theorem]\label{prob:vii11-07}
Translate Stokes for $\iota_XdV$ into flux equals volume divergence.
\end{problem}
\begin{solution}
The identity $d(\iota_XdV)=(\operatorname{div}X)dV$ supplies the interior term. Restricting $\iota_XdV$ to the oriented boundary evaluates the normal component of $X$ times surface measure.
\end{solution}

\begin{problem}[Annulus signs]\label{prob:vii11-08}
Determine the induced orientations on the two boundary circles of a standard oriented annulus.
\end{problem}
\begin{solution}
Outward normal on the outer circle points radially outward, giving counterclockwise tangent. On the inner circle outward relative to the annulus points radially inward, forcing the positive tangent to be clockwise.
\end{solution}

\begin{problem}[Exact area-form obstruction]\label{prob:vii11-09}
Show the standard area form on $S^2$ is not exact.
\end{problem}
\begin{solution}
If $\omega=d\beta$, Stokes on boundaryless $S^2$ gives $\int_{S^2}\omega=0$, contradicting the positive area $4\pi$.
\end{solution}

\begin{problem}[Zero circulation of gradients]\label{prob:vii11-10}
Show $\int_Cdf=0$ for every closed oriented curve $C$.
\end{problem}
\begin{solution}
The one-form $df$ is exact. Applying Stokes to the closed one-manifold $C$ gives zero; equivalently parameter endpoints cancel.
\end{solution}

\begin{problem}[Flux of radial field]\label{prob:vii11-11}
Compute the flux of $X(x)=x$ across the unit sphere in $\mathbb R^n$ using divergence.
\end{problem}
\begin{solution}
$\operatorname{div}X=n$. Thus the flux is $n\operatorname{Vol}(B^n)$, equal to the $(n-1)$-area of $S^{n-1}$.
\end{solution}

\begin{problem}[Boundary of boundary cancellation]\label{prob:vii11-12}
Explain the sign cancellation at corners of an oriented rectangle.
\end{problem}
\begin{solution}
Each vertex is reached from two oriented edges with opposite endpoint signs. Hence the signed $0$-chain $\partial^2R$ vanishes, the geometric analogue of $d^2=0$.
\end{solution}


\section{Challenge problems}

\begin{problem}[Stokes and homotopy invariance preview]\label{prob:vii11-ch01}
Let $C_0$ and $C_1$ be oriented closed curves bounding an oriented surface strip $S$. If $\alpha$ is closed, show $\int_{C_1}\alpha=\int_{C_0}\alpha$.
\end{problem}
\begin{solution}
The oriented boundary is $C_1-C_0$. Therefore $\int_{C_1}\alpha-\int_{C_0}\alpha=\int_{\partial S}\alpha=\int_Sd\alpha=0$. This is the basic mechanism behind homotopy invariance of periods.
\end{solution}

\begin{problem}[No compact manifold with one regular coordinate function]\label{prob:vii11-ch02}
Suppose a compact connected oriented boundaryless manifold $M$ admitted an $(n-1)$-form $\eta$ with $d\eta$ everywhere positive. Derive a contradiction.
\end{problem}
\begin{solution}
Positivity gives $\int_Md\eta>0$. But Stokes on a boundaryless compact manifold gives $\int_Md\eta=0$.
\end{solution}

\begin{problem}[Gauss-law form]\label{prob:vii11-ch03}
On $\mathbb R^n\setminus\{0\}$ construct the standard radial flux $(n-1)$-form and explain why its integrals over concentric spheres agree.
\end{problem}
\begin{solution}
Take $\omega=\iota_XdV/\|x\|^n$ with $X=\sum x^i\partial_i$. A direct computation gives $d\omega=0$ away from the origin. Apply Stokes to the annulus between two spheres; their induced orientations yield equality of the two flux integrals.
\end{solution}

\begin{problem}[Exactness detected by an integral]\label{prob:vii11-ch04}
Use Stokes to prove that the form $\alpha=(-y\,dx+x\,dy)/(x^2+y^2)$ on $\mathbb R^2\setminus\{0\}$ is closed but not exact.
\end{problem}
\begin{solution}
A direct calculation gives $d\alpha=0$. On the unit circle, $\alpha=d\theta$ along the parametrization and integrates to $2\pi$. Every exact one-form has zero integral on a closed curve, so $\alpha$ cannot be exact.
\end{solution}

\begin{problem}[Boundary orientation under products]\label{prob:vii11-ch05}
For oriented manifolds with boundary, derive the sign formula $\partial(M^m\times N)=\partial M\times N\;\cup\;(-1)^m M\times\partial N$.
\end{problem}
\begin{solution}
Compare outward-normal-first frames with the product frame. For the second boundary piece, the outward normal from $N$ must be moved past the $m$ tangent vectors of $M$, producing the factor $(-1)^m$.
\end{solution}

\section*{Bridge to VII/12}
Stokes' theorem completes the first differential-form arc. VII/12 returns to one-dimensional geometry: regular parametrized curves, reparametrization, arc length, unit tangent, and line integrals, preparing the Frenet invariants of VII/13.
